\documentclass[12pt]{article}
\usepackage[utf8]{inputenc}
\usepackage[T1]{fontenc}
\usepackage{hyperref}
\title{Tulu Language Resources for K-12 Curriculum Development}
\author{Compiled by Claude Sonnet 5 (claude-sonnet-5) for Andr\'as Kornai}
\date{}
\begin{document}
\maketitle

\section{Introduction}
Tulu (ISO 639-3 tcy; Glottocode tulu1258) is a Dravidian language of coastal Karnataka and northern Kerala, India, with about 2,000,000 speakers. The reference description says it is not known to be taught in schools. Search results indicate otherwise: the Karnataka Tulu Sahitya Academy is reported to have introduced Tulu and the Tigalari-Tulu script into schools in the Mangalore and Udupi districts. This needs verification. The sources are mostly secondary and some are dated.

\section{Primary Educational Resources}
\begin{itemize}
\item ``A handbook to learn Tulu'' by Pradyoth Hegde (2020), listed in a Unicode submission by KaNaada Phonetics. I could not confirm its availability: \url{https://www.unicode.org/L2/L2022/22075-tulu-response-letter.pdf}
\end{itemize}

\section{Linguistic and Reference Materials}
\begin{itemize}
\item U.\,P. Upadhyaya, \emph{Concise Tulu Dictionary} (5 vols, 2005), cited in the same Unicode document: \url{https://www.unicode.org/L2/L2022/22075-tulu-response-letter.pdf}
\item Online Tulu dictionary at tuluworld.org (English and Kannada meanings), described in a 2012 blog post. Current availability is unverified: \url{https://www.raveeshkumar.com/2012/02/tulu-language-resources-online.html}
\item Wikipedia article on the Tulu language: \url{https://en.wikipedia.org/wiki/Tulu_language}
\item Omniglot page on the Tulu writing system: \url{https://omniglot.com/writing/tulu.htm}
\item Omniglot page on the Tigalari script: \url{https://www.omniglot.com:443/writing/tigalari.htm}
\item Endangered Alphabets page on the Tulu script: \url{https://endangeredalphabets.net/?p=438}
\end{itemize}

\section{Religious and Community Resources}
\begin{itemize}
\item Karnataka Tulu Sahitya Academy (founded 1994) and the Kerala Tulu Academy, named as regulating bodies on Wikipedia: \url{https://en.wikipedia.org/wiki/Tulu_language}
\end{itemize}

\section{Digital and Audio Resources}
\begin{itemize}
\item The Unicode submission mentions Tulu Unicode fonts and an Android keyboard, Kanatavu (2020): \url{https://www.unicode.org/L2/L2022/22075-tulu-response-letter.pdf}
\end{itemize}

\section{Curriculum Development Context}
Tulu has a large speaker base, an active academy and a script revival (Tigalari-Tulu) with Unicode support in progress. Reported school introduction in Mangalore and Udupi should be verified with the Academy. Brigel's 1872 grammar was mentioned in search results via Open Library, but I did not see a direct URL, so it is not listed.
\end{document}